% s-307-p307-m002; original PDF 307, 308
\Aphorism{s-307-p307-m002}{361}
Bei welcher Art Menschen findet der Pessimismus Gehör, jene wunderliche Philosophie, die dem Schauspiel „Welt“ eine Kritik schuldig zu sein glaubt und diese Kritik auf das \Sperr{Leiden} in der Welt beschränkt, von der simplen langen Weile bis hinauf zu den Siedegraden des Schmerzes? Gewiss nicht bei Solchen, die selbst eine starke Leidenserfahrung und damit die natürliche Furcht und Feigheit vorm Leiden hinter sich haben; kaum wohl bei Denen, die zu stumpf und zäh sind, um zu leiden oder sich Leiden vorzustellen, sondern — nun, bei den Übrigbleibenden, bei Solchen, die hin und wieder ein mittelmässiger Schmerz kitzelt, stark genug, um ihre Phantasie zu schrecklichen Verallgemeinerungen und Al-fresco-Malereien anzureizen, aber nicht stark genug, um abzuhärten und das Phantasma in der Wirklichkeit zu ersticken. Wer selbst hohe Foltergrade überstanden hat, wird über die ganze Welt der Folterkammern, über das Kreuz auf Golgatha und den Stier des Phalaris, über die Schrecknisse des Naturverlaufs und der Weltgeschichte etwas geringschätzig denken; umgekehrt wird so ein halbgeplagter Zärtling nicht ermangeln, seine Kolik oder seinen Zahnschmerz in gerader Linie ins Metaphysische hinein zu projiciren, als Rathschluss und Heimsuchung Gottes, Vorgeschmack der Hölle, Schuld des Willens zum Leben, \Greek{δεύτερος} \Greek{πλοῦς} zur Erlösung, oder je nach dem Lexicon seiner Philosophie. Für Jenen haben die Dinge einen helleren, friedlicheren Anblick gewonnen, für Diesen wirkt die imaginäre Verlängerung unvollständiger Leidensgrade in einer Weise verdüsternd, die bis zum thatsächlichen Martyrium führen kann. Das fortdauernde Vergegenwärtigen und Mitempfinden fremden und vergangenen Leidens, namentlich die Betrachtung eines bestimmten extremen Falls, noch unterstützt durch ein gewisses, nicht sehr reiches Maass eigener Leidenserfahrung, lässt einen hypochondrischen Zweifel bestehen, ob derlei Extreme überhaupt menschenmöglich sind — man wird zum Märtyrer, um sich diesen Zweifel zu lösen. Aus der Ungewissheit, wieviel wir maximaliter an Schmerz vertragen können, verlangen wir, in einer Anwandlung verzweifelter Glücksspielerei, nach Kreuz und Rost, nach einer Tiefe des Leidens, hinter der nicht noch unerschöpfte Tiefen qualvolleren Leidens als bedrohliche Möglichkeiten offen bleiben.
\par

